\caption{\textbf{The cortical--subcortical divide in response types is mainly a firing-rate divide.}
All trials, 47,378 neurons in 515 insertions (the neurons of Fig.~3a with a cortical or subcortical Cosmos label). Cortical: Isocortex, OLF, HPF, CTXsp (green in the Swanson maps); subcortical: CNU, TH, HY, MB, HB, CB (pink/orange); neurons with Beryl or Cosmos label root or void removed. Firing rate: mean spike count of the feature vector (all PETH windows), in Hz.
\textbf{a} Firing rate per Cosmos area (violins; dots, medians; bars, interquartile range). Subcortical neurons fire about twice as fast as cortical neurons (median 6.9 vs 3.5\,Hz).
\textbf{b} Median rate of each Beryl region ($\geq$20 neurons; 198 regions, 68 cortical), ranked, coloured by Cosmos area as in \textbf{a}. Inset: fraction of cortical regions in a running window of 21 ranks (dotted, overall fraction). Regional median rate separates cortical from subcortical regions (AUC 0.87); it does not depend on the number of neurons recorded per region (Fig.~5n).
\textbf{c} Decoding a single neuron's cortical or subcortical location (logistic regression, 5-fold cross-validation grouped by insertion, so test probes are never seen in training; AUC) from its rate, odd/even reliability, peak stimulus latency (as Fig.~5f), k-means cluster (Fig.~3a), response shape (30 principal components of the $z$-scored feature vector) or shape and rate. Bottom: rate-matched neurons (cortical and subcortical neurons subsampled to the same log-rate histogram; mean of 5 resamples, 34,152 neurons each). Rate is the best single predictor (0.64); response shape adds little to it (0.66); on rate-matched neurons, response type and shape keep a small separation (0.58).
\textbf{d} The cortical--subcortical separation of old Fig.~S11, recomputed: cell-weighted Fisher criterion $D$ between the regions' 25-cluster composition vectors (fraction of the region's neurons in each k-means cluster; regions $\geq$20 neurons), dots; grey, 95\% range of region-label shuffles (1,000 for all neurons, 200 per resample otherwise; line, median). Response-type composition, all neurons: $D=0.62$ (shuffle $0.20\pm0.05$, $p=0.001$). Random subsamples of the rate-matched size (20 resamples): $D=0.60$. Rate-matched (20 resamples): $D=0.50$, still above all shuffles ($p=0.005$). The same statistic on 25 rate-quantile bins instead of response types: $D=1.01$: regions differ more in their rate distributions than in their response types.
\textbf{e} k-means clusters of Fig.~3a (numbers): fraction of cortical neurons vs median rate (dot size, cluster size; dotted, overall cortical fraction). Low-rate clusters are more cortical (Spearman $\rho=-0.39$).
\textbf{f} As \textbf{e} for the 100 Rastermap clusters of Fig.~5a (row spans, as there; $\rho=-0.63$); blue, sparse-tiling clusters; orange, dense-tiling clusters.
\textbf{g} Sparse-tiling neurons and dense tiling (Fig.~5) as a percentage of cortical (dark) and subcortical (light) neurons, all neurons and rate-matched (20 resamples). Sparse-tiling neurons are 1.64 times more frequent in cortex (28.4\% vs 17.3\%), 1.29 times after rate matching; dense tiling is about equally frequent (0.86, rate-matched 1.04).}
\label{fig:cortex_rate}
